\documentclass[11pt]{article}
\usepackage[spanish,es-tabla,es-nodecimaldot,es-noshorthands]{babel}
\usepackage[utf8]{inputenc}
\usepackage[T1]{fontenc}
\usepackage{lmodern}
\usepackage[margin=2.5cm]{geometry}
\usepackage{amsmath}
\usepackage{booktabs}
\usepackage{tabularx}
\usepackage{float}
\usepackage{enumitem}
\usepackage{xcolor}
\usepackage{listings}
\usepackage{tikz}
\usetikzlibrary{arrows.meta, positioning, calc, fit, decorations.pathreplacing, babel}
\usepackage[hidelinks]{hyperref}

\setlist{itemsep=2pt, topsep=4pt}
\newcolumntype{Y}{>{\raggedright\arraybackslash}X}

\lstset{
  language=C++, basicstyle=\footnotesize\ttfamily, columns=fullflexible,
  keepspaces=true, frame=lines, framesep=4pt, xleftmargin=4pt,
  xrightmargin=4pt, commentstyle=\color{black!55}\itshape,
  keywordstyle=\bfseries, showstringspaces=false, aboveskip=8pt,
  belowskip=6pt, captionpos=b,
}
\renewcommand{\lstlistingname}{Listado}

\tikzset{
  box/.style   = {draw=black!75, semithick, rounded corners=1.5pt, align=center,
                  minimum height=1.1cm, minimum width=2.6cm, font=\footnotesize,
                  fill=white},
  shade/.style = {box, fill=black!7},
  sig/.style   = {-{Stealth[length=2mm]}, semithick},
  lab/.style   = {font=\scriptsize, midway, above},
  note/.style  = {font=\scriptsize\itshape, align=center},
  panel/.style = {draw=black!40, rounded corners=3pt, inner sep=8pt},
}

\newcommand{\code}[1]{\texttt{#1}}
\newcommand{\us}{\,$\mu$s}

\title{El mutex de los bloques de Simulink\\[2pt]
       \large Por qué no causaba el bloqueo del enlace USB--SPI\\
       y qué se cambió para eliminarlo}
\author{Haniel Ulises\\Servicio social}
\date{11 de octubre de 2026}

\begin{document}
\maketitle

\begin{abstract}
Se responde la pregunta de si el mutex de los bloques de Simulink
\code{escribirPic} y \code{leerPic} es compatible con Windows y si es la causa
del bloqueo de la comunicación con el dsPIC33CK64MC105. El mutex es un objeto
del núcleo de Windows (\code{CreateMutex}, \code{WaitForSingleObject},
\code{ReleaseMutex}) y es compatible con ese sistema, pero no puede causar el
bloqueo: Simulink ejecuta los bloques uno después del otro en un solo hilo, por
lo que el mutex siempre está libre cuando un bloque lo pide. El bloqueo y las
tramas corruptas tenían su origen en el firmware del dsPIC, que podía quedarse
en una espera activa y dejaba de recibir al desbordarse la FIFO del SPI, y en
el uso que la PC hacía del bus, con una escritura y una lectura separadas en
cada paso. Se describen los cambios realizados para eliminar esas causas: un
protocolo de una sola transferencia full-duplex de ocho bytes por paso,
delimitada por \code{CS} y protegida con número de secuencia y CRC-8; un
firmware nuevo que atiende el SPI con DMA; bloques de Simulink en los que un
solo bloque usa el FT2232H, sin mutex; y herramientas para probar el enlace sin
MATLAB. En Windows, el modelo \code{MergedPic.slx} del asesor, sin
modificaciones y con los bloques nuevos, ejecutó 10\,001 pasos a 1\,kHz sin
errores de comunicación.
\end{abstract}

\tableofcontents
\newpage

% ===========================================================================
\section{La pregunta y la respuesta corta}

La pregunta fue si el mutex que usan los bloques de Simulink es compatible con
Windows y si el bloqueo de la comunicación con el dsPIC se debía a él.

\begin{itemize}
\item \textbf{¿Es compatible con Windows?} Sí. No es un mecanismo de Simulink
  ni de C++, sino un objeto del núcleo de Windows creado y usado con las
  funciones de \code{kernel32}. Compila con MSVC y funciona en cualquier
  versión de Windows; de hecho, sólo existe en Windows.
\item \textbf{¿Es la causa del bloqueo?} No. En una simulación de Simulink los
  bloques se ejecutan en un solo hilo, de modo que nunca hay dos bloques
  esperando el mutex al mismo tiempo. Además, un mutex de Windows no bloquea al
  hilo que ya lo tiene, y si su dueño desaparece el siguiente que lo pide
  continúa. La sección~\ref{sec:mutex} detalla los cuatro argumentos.
\item \textbf{¿De dónde venía el bloqueo?} Del firmware del dsPIC y del uso del
  bus SPI desde la PC (sección~\ref{sec:causas}). Esas causas ya estaban
  identificadas en la propuesta (\code{docs/propuesta.pdf}) y en la sección
  ``Limitaciones conocidas'' del \code{README.md}.
\end{itemize}

Los bloques nuevos no tienen mutex porque un solo bloque usa el FT2232H, pero
eso es una consecuencia del rediseño, no la corrección del bloqueo.

% ===========================================================================
\section{Cómo funcionaban los bloques anteriores}
\label{sec:anteriores}

Los bloques anteriores eran tres S-Functions de nivel~2 que compartían el
FT2232H. El código está en \code{SIMULINK/WritePicModificado} y
\code{SIMULINK/ReadPicModificado}, y el modelo que los usa es
\code{MergedPic.slx}. La tabla~\ref{tab:anteriores} resume lo que hacía cada
uno.

\begin{table}[H]
\centering
\small
\begin{tabularx}{\linewidth}{@{}lYl@{}}
\toprule
Bloque & Qué hacía & Operación SPI por paso \\
\midrule
\code{initializePic} & Abría el canal~0 de libMPSSE a 1\,MHz en modo~0, con
  \code{CS} activo en bajo en \code{ADBUS3}, y creaba el mutex. Su salida era
  la ``llave'': la dirección de la estructura compartida \code{HandleFTDI}
  expresada como entero de 64 bits. & Ninguna \\
\code{escribirPic} & Convertía el voltaje a un código de 16 bits,
  $13107\,(V + 2.5)$. Si el código no había cambiado respecto al paso anterior,
  no enviaba nada. & \code{SPI\_Write}, 3 bytes \\
\code{leerPic} & Leía 12 bytes y buscaba la secuencia \code{0xAA 0x55}
  seguida de 4 bytes de posición. & \code{SPI\_Read}, 12 bytes \\
\bottomrule
\end{tabularx}
\caption{Bloques anteriores.}
\label{tab:anteriores}
\end{table}

La estructura compartida está definida en \code{ftdi\_compartido.h}
(listado~\ref{lst:estructura}). Contiene el \emph{handle} del FT2232H, una
bandera que indica si el canal se abrió y el \emph{handle} del mutex. El
\code{\#include <mutex>} de la biblioteca estándar de C++ no se usa: el mutex
es el de Windows.

\noindent\begin{minipage}{\linewidth}
\begin{lstlisting}[caption={Estructura compartida
  (\code{SIMULINK/WritePicModificado/ftdi\_compartido.h}).},
  label=lst:estructura]
#include <mutex>        // no se usa
#include <windows.h>
#include "ftd2xx.h"
#include "libmpsse_spi.h"

struct HandleFTDI {
    FT_HANDLE handle;
    bool abierto;
    HANDLE semaforo;    // mutex del nucleo de Windows
};
\end{lstlisting}
\end{minipage}

\code{initializePic} crea el mutex sin nombre en \code{mdlStart}
(\code{initializePic.cpp}, línea~52) y lo destruye con \code{CloseHandle} en
\code{mdlTerminate}, después de enviar \code{0xAA 0x00 0x00}. Con la
conversión de \code{escribirPic}, el código \code{0x0000} equivale a
$-2.5$\,V, no a 0\,V. \code{escribirPic} y \code{leerPic} toman el mutex antes
de su operación sobre el bus y lo sueltan al terminar
(listado~\ref{lst:uso}).

\noindent\begin{minipage}{\linewidth}
\begin{lstlisting}[caption={Uso del mutex en \code{escribirPic.cpp}
  (líneas 174--220) y \code{leerPic.cpp} (líneas 110--125), resumido.},
  label=lst:uso]
// escribirPic, mdlOutputs
if (!primerEnvio[0] && valor_dac == ultimoValor[0])
    return;                                   // valor repetido: no envia
WaitForSingleObject(hw->semaforo, INFINITE);
bufferTx[0] = 0xAA;
bufferTx[1] = (valor_dac >> 8) & 0xFF;
bufferTx[2] = valor_dac & 0xFF;
SPI_Write(hw->handle, bufferTx, 3, &bytesTransferidos, ...);
ReleaseMutex(hw->semaforo);

// leerPic, mdlOutputs
WaitForSingleObject(hw->semaforo, INFINITE);
status = SPI_Read(hw->handle, Buf, 12, &bytesTransferidos, ...);
ReleaseMutex(hw->semaforo);
for (i = 0; i < bytesTransferidos - 5; i++)
    if (Buf[i] == 0xAA && Buf[i+1] == 0x55) { /* posicion en Buf[i+2..i+5] */ }
\end{lstlisting}
\end{minipage}

En cada paso de simulación había entonces dos transacciones independientes
sobre el bus: una escritura de 3 bytes y una lectura de 12. La parte izquierda
de la figura~\ref{fig:antes-ahora} lo muestra.

% ===========================================================================
\section{El mutex}
\label{sec:mutex}

\subsection{Es compatible con Windows}

Las cuatro funciones que usan los bloques, \code{CreateMutex},
\code{WaitForSingleObject}, \code{ReleaseMutex} y \code{CloseHandle}, son del
API de Windows y están en \code{kernel32.dll}, que está presente en todas las
versiones. El mutex se crea sin nombre, por lo que sólo se puede alcanzar a
través del \emph{handle} guardado en \code{HandleFTDI}: ningún otro proceso
puede tomarlo. No hay ninguna incompatibilidad entre este mutex y Windows,
MATLAB o el compilador MSVC.

\subsection{No puede causar el bloqueo}

Un mutex sólo detiene a un hilo cuando otro hilo distinto lo tiene tomado y no
lo suelta. En estos bloques esa situación no se presenta, por cuatro razones.

\begin{enumerate}
\item \textbf{Nunca hay dos bloques compitiendo por él.} En una simulación
  normal, Simulink calcula las salidas de los bloques en un orden fijo y en un
  solo hilo: llama a \code{mdlOutputs} de \code{escribirPic}, regresa, y
  después llama a \code{mdlOutputs} de \code{leerPic}
  (figura~\ref{fig:hilo}). Cuando un bloque pide el mutex, el otro ya lo soltó,
  y \code{WaitForSingleObject} regresa de inmediato. El mutex nunca espera.
\item \textbf{Un mutex de Windows no bloquea al hilo que ya lo tiene.} Es
  recursivo para su dueño: el mismo hilo puede volver a tomarlo sin quedarse
  esperando. Aunque faltara una llamada a \code{ReleaseMutex} en algún camino
  del código, el hilo de la simulación no se detendría por ello.
\item \textbf{Tampoco espera si su dueño desaparece.} Si el hilo que tiene el
  mutex termina sin soltarlo, Windows lo marca como abandonado y el siguiente
  \code{WaitForSingleObject} regresa con \code{WAIT\_ABANDONED} en lugar de
  esperar para siempre.
\item \textbf{No hay otro candado en la cadena.} Se revisaron las funciones que
  importa \code{libmpsse.dll}: de \code{kernel32} sólo usa \code{Sleep},
  \code{QueryPerformanceCounter}, \code{LoadLibraryW}, \code{GetProcAddress} y
  funciones de diagnóstico, y carga \code{ftd2xx.dll} en tiempo de ejecución.
  No importa ninguna primitiva de sincronización (mutex, sección crítica ni
  evento). Las lecturas y escrituras del D2XX tienen tiempos límite
  (\code{FT\_SetTimeouts}), por lo que tampoco esperan indefinidamente.
\end{enumerate}

\begin{figure}[H]
\centering
\begin{tikzpicture}[x=1cm, y=1cm,
    seg/.style = {draw=black!75, semithick, minimum height=0.75cm,
                  font=\scriptsize, inner sep=1pt, anchor=west},
    mtx/.style = {seg, fill=black!7}]
  \node[font=\footnotesize, anchor=east] at (-0.2,0) {Hilo de Simulink};
  \draw[->, semithick, black!60] (0,-0.6) -- (15.4,-0.6)
    node[right, font=\scriptsize] {$t$};
  % escribirPic
  \node[mtx, minimum width=1.0cm] (w1) at (0,0) {Wait};
  \node[seg, minimum width=2.4cm] (w2) at (w1.east) {\code{SPI\_Write} 3\,B};
  \node[mtx, minimum width=1.2cm] (w3) at (w2.east) {Release};
  % leerPic
  \node[mtx, minimum width=1.0cm] (r1) at ($(w3.east)+(0.3,0)$) {Wait};
  \node[seg, minimum width=2.4cm] (r2) at (r1.east) {\code{SPI\_Read} 12\,B};
  \node[mtx, minimum width=1.2cm] (r3) at (r2.east) {Release};
  % paso siguiente
  \node[mtx, minimum width=1.0cm] (n1) at ($(r3.east)+(1.4,0)$) {Wait};
  \node[seg, minimum width=2.4cm, draw=black!40, text=black!50] (n2)
    at (n1.east) {\code{SPI\_Write}\ldots};
  % llaves
  \draw[decorate, decoration={brace, amplitude=4pt}, semithick]
    (w1.north west) ++(0,0.08) -- ($(w3.north east)+(0,0.08)$)
    node[midway, above=4pt, font=\scriptsize] {\code{escribirPic}: \code{mdlOutputs}};
  \draw[decorate, decoration={brace, amplitude=4pt}, semithick]
    (r1.north west) ++(0,0.08) -- ($(r3.north east)+(0,0.08)$)
    node[midway, above=4pt, font=\scriptsize] {\code{leerPic}: \code{mdlOutputs}};
  \node[font=\scriptsize, above=6pt] at ($(n1.north)!0.5!(n2.north)$) {paso $k+1$};
  % notas
  \node[note, below=0.55cm] at (w1.south) {libre:\\regresa de inmediato};
  \node[note, below=0.55cm] at (r1.south) {libre:\\regresa de inmediato};
\end{tikzpicture}
\caption{Un paso de simulación con los bloques anteriores. Los dos bloques se
  ejecutan uno tras otro en el mismo hilo; cada \code{Wait} encuentra el mutex
  libre.}
\label{fig:hilo}
\end{figure}

\subsection{Qué hacía realmente}

El mutex sólo evitaba que dos hilos usaran el FT2232H al mismo tiempo, algo que
en este modelo no ocurre. Podría ser útil si los bloques se ejecutaran en
hilos distintos, por ejemplo en un modelo con varias tareas concurrentes, pero
no es el caso de \code{MergedPic.slx}. No interviene en ninguna de las fallas
de la sección~\ref{sec:causas}: con o sin mutex, la PC hacía las mismas dos
transacciones en el mismo orden.

% ===========================================================================
\section{Las causas reales del bloqueo}
\label{sec:causas}

El bloqueo venía del firmware anterior del dsPIC (\code{dsPicController}), que
podía quedarse detenido o dejar de recibir, y de la forma en que la PC usaba el
bus. La tabla~\ref{tab:causas} reúne las causas.

\begin{table}[H]
\centering
\small
\begin{tabularx}{\linewidth}{@{}lYY@{}}
\toprule
Dónde & Problema & Efecto \\
\midrule
Firmware & \textbf{Espera activa.} La escritura en SPI1 esperaba mientras la
  FIFO de transmisión estuviera llena, condición que sólo cambia cuando la PC
  genera reloj. & El lazo principal se detenía y no atendía la recepción. \\
\addlinespace
Firmware & \textbf{Desbordamiento sin recuperación.} Con la FIFO de recepción
  llena se activa \code{SPIROV}, y el driver generado por MCC no limpiaba la
  bandera. & SPI1 dejaba de recibir para siempre: el bloqueo permanente. \\
\addlinespace
Firmware & \textbf{Respuesta sin sincronía.} El esclavo escribía la posición
  en la FIFO cada vez que había espacio, sin relación con el momento en que la
  PC leía. & Datos desfasados y desalineados. \\
\addlinespace
Firmware & \textbf{Sin delimitación de tramas.} No usaba \code{CS}; sólo
  buscaba \code{0xAA}, valor que puede aparecer dentro de los datos. & Tramas
  mal interpretadas. \\
\addlinespace
PC & \textbf{Bus full-duplex usado como half-duplex.} \code{SPI\_Write} y
  \code{SPI\_Read} por separado. Al escribir se descartaba lo que llegaba por
  MISO; al leer, el dsPIC tomaba los bytes de relleno como datos. & Dos
  transacciones por paso que se corrompían entre sí y llenaban la FIFO del
  dsPIC. \\
\addlinespace
PC & \textbf{\code{escribirPic} no reenviaba un valor repetido.} & Si una
  escritura se perdía, con una entrada constante el error era permanente. \\
\addlinespace
PC & \textbf{Al terminar se enviaba \code{0x0000}.} & $-2.5$\,V en lugar de
  0\,V: el motor quedaba energizado. \\
\bottomrule
\end{tabularx}
\caption{Causas del bloqueo y de las tramas corruptas.}
\label{tab:causas}
\end{table}

La combinación de las primeras dos explica el bloqueo. En el firmware
anterior, el lazo principal escribía los 6 bytes de la respuesta con
\code{SPI1\_ByteWrite}, que espera mientras \code{SPITBF} indique la FIFO de
transmisión llena. Durante esa espera nadie leía la FIFO de recepción de 4
niveles, a la que llegaban los 3 bytes de cada escritura y los 12 de relleno de
cada lectura de la PC. Al llenarse se activaba \code{SPIROV} y, como ninguna
parte del código lo borraba, SPI1 ya no recibía nada más. Ninguno de estos
pasos depende del mutex.

En la primera prueba en Windows (sección~\ref{sec:resultados}) apareció otra
causa, ajena al firmware: en esa PC libMPSSE numeró la interfaz~B del FT2232H
como canal~0. Los bloques abrían el canal~0, así que hablaban por la interfaz
que no está cableada al dsPIC y sólo leían \code{0xFF}. Los bloques nuevos
buscan la interfaz por la que responde el dsPIC (sección~\ref{sec:bloques}),
de modo que ya no dependen de esa numeración.

% ===========================================================================
\section{El protocolo nuevo}
\label{sec:protocolo}

Cada paso de simulación hace una sola transferencia full-duplex de 8 bytes con
\code{SPI\_ReadWrite}, con \code{CS} en bajo durante toda la trama. Mientras la
PC envía el voltaje por MOSI, el dsPIC devuelve la posición por MISO en los
mismos ocho bytes de reloj. El formato está definido en un solo archivo,
\code{protocolo/daq\_protocolo.h}, que comparten el firmware y la PC
(tabla~\ref{tab:trama}).

\begin{table}[H]
\centering
\small
\begin{tabularx}{\linewidth}{@{}cYY@{}}
\toprule
Byte & PC $\to$ dsPIC (MOSI) & dsPIC $\to$ PC (MISO) \\
\midrule
0 & Encabezado \code{0xA5} & Encabezado \code{0x5A} \\
1 & Número de secuencia & Eco de la secuencia de la última trama válida \\
2 & Código del DAC, byte alto & Posición del encoder, \code{int32} \\
3 & Código del DAC, byte bajo & little-endian (bytes 2 a 5) \\
4 & Banderas: bit~0 salida habilitada, bit~1 reinicio del encoder & \\
5 & Reservado (0) & \\
6 & Reservado (0) & Status: errores de CRC, encabezado, longitud y
  vigilancia; causa de reinicio en los 4 bits altos \\
7 & CRC-8 de los bytes 0 a 6 & CRC-8 de los bytes 0 a 6 \\
\bottomrule
\end{tabularx}
\caption{Formato de la trama de 8 bytes.}
\label{tab:trama}
\end{table}

Qué resuelve cada parte:

\begin{itemize}
\item \textbf{Una transferencia por paso.} Ya no hay una escritura y una
  lectura separadas que se estorben: cada byte que entra al dsPIC es parte de
  una trama válida y cada byte que sale es parte de la respuesta.
\item \textbf{\code{CS} delimita la trama.} El dsPIC sabe dónde empieza y
  termina sin buscar un valor especial en los datos.
\item \textbf{Secuencia y CRC-8} (polinomio \code{0x07}, CRC-8/SMBUS). La PC
  detecta tramas dañadas, perdidas o repetidas, y las cuenta como errores.
\item \textbf{El voltaje se envía en cada paso}, aunque no cambie, de modo que
  una trama dañada sólo afecta un periodo.
\item \textbf{Bandera de salida habilitada.} Con ella en 0, el dsPIC lleva el
  DAC a 0\,V (código 32767). Al terminar la simulación se envía una trama así,
  en lugar de $-2.5$\,V.
\item \textbf{Status.} El dsPIC informa sus propios errores y, una vez después
  de cada arranque, la causa del reinicio (encendido, MCLR, \emph{watchdog} o
  tipo de trampa).
\end{itemize}

La respuesta de la transferencia $k$ se arma al terminar la $k-1$, por lo que
la posición llega con un paso de retraso. Sumado al paso en que se aplica el
voltaje, el retardo total del lazo es de dos periodos de muestreo.

% ===========================================================================
\section{El firmware del dsPIC: \code{dsPicDaq}}
\label{sec:firmware}

El firmware nuevo ya no tiene esperas activas sobre el SPI y recupera el
desbordamiento en cada trama, que eran las dos causas del bloqueo. El proyecto
está en \code{MPLAB/dsPicDaq.X} y la imagen compilada en
\code{MPLAB/dsPicDaq.hex}.

\begin{itemize}
\item \textbf{DMA en lugar de espera activa.} SPI1 trabaja como esclavo en
  modo de buffer estándar. DMA0 recibe y DMA1 transmite durante la
  transferencia, sin intervención del lazo principal. La FIFO del modo mejorado
  no se usa: tiene 4 niveles de 8 bits, menos que una trama.
\item \textbf{Fin de trama por \code{CS}.} \code{CS} (RB10) se asigna con PPS a
  \code{SS1} y a \code{INT1}. En el flanco de subida, la interrupción valida la
  trama, reinicia SPI1 (vacía los buffers y borra \code{SPIROV}), muestrea el
  QEI, arma la respuesta siguiente y reprograma el DMA. Aunque una trama llegue
  incompleta o dañada, la siguiente empieza con el módulo limpio.
\item \textbf{DAC.} El lazo principal sólo escribe al DAC8554 cuando el valor
  cambia.
\item \textbf{Vigilancia.} Timer1: si pasan 50\,ms sin una trama válida, el DAC
  va a 0\,V y se reporta en el status. Si Simulink se detiene, el motor no se
  queda con el último voltaje.
\item \textbf{Reloj.} El PLL eleva la frecuencia de instrucción de 4\,MHz a
  $F_{cy} = 50$\,MHz. Con 4\,MHz, la interrupción de fin de trama tardaba cerca
  de 100\us{} y el esclavo no alcanzaba a presentar a tiempo el primer bit de
  cada byte a 1\,MHz.
\item \textbf{Causa de reinicio.} \code{reinicio.c} guarda el código de una
  trampa en RAM persistente antes del reinicio, y el dsPIC lo reporta en la
  primera respuesta después de arrancar.
\item \textbf{Medición.} LED0 (RD10) queda en alto durante la interrupción,
  para medir su duración con un analizador lógico.
\end{itemize}

El firmware se compila con \code{-O1} para acortar la interrupción de fin de
trama, que fija el intervalo mínimo entre tramas. En Linux se compila sin MPLAB
X con \code{MPLAB/dsPicDaq.X/compilar.sh}.

% ===========================================================================
\section{Los bloques nuevos de Simulink}
\label{sec:bloques}

Hay tres juegos de bloques en \code{SIMULINK/DaqPic}, y en todos un solo bloque
usa el FT2232H, una vez por paso. Por eso ya no hace falta el mutex. Se
compilan todos con \code{compilar.m} (tabla~\ref{tab:bloques}).

\begin{table}[H]
\centering
\small
\begin{tabularx}{\linewidth}{@{}lYY@{}}
\toprule
Juego & Bloques & Uso \\
\midrule
Un bloque & \code{daqPic}: entra el voltaje; salen posición, errores y
  atrasos & Modelos nuevos \\
\addlinespace
Tres bloques & \code{daqPicInicio} (llave), \code{daqPicEscribir} (llave,
  voltaje), \code{daqPicLeer} (posición, errores, atrasos) & Modelos que
  separan inicialización, escritura y lectura \\
\addlinespace
Nombres anteriores & \code{initializePic}, \code{escribirPic},
  \code{leerPic}, con los mismos puertos y en el mismo orden & Usar
  \code{MergedPic.slx}, \code{ReadPicModificado.slx} y
  \code{WritePicModificado.slx} sin modificarlos \\
\bottomrule
\end{tabularx}
\caption{Bloques nuevos.}
\label{tab:bloques}
\end{table}

\begin{figure}[H]
\centering
\begin{tikzpicture}[node distance=0.7cm and 0.9cm]
  % ---------------- Antes ----------------
  \node[box] (ai) {\code{initializePic}\\[-1pt]\scriptsize abre el canal 0,\\[-2pt]
    \scriptsize crea el mutex};
  \node[box, below left=1.0cm and -0.9cm of ai] (aw)
    {\code{escribirPic}\\[-1pt]\scriptsize mutex +\\[-2pt]
     \scriptsize \code{SPI\_Write} 3\,B};
  \node[box, below right=1.0cm and -0.9cm of ai] (ar)
    {\code{leerPic}\\[-1pt]\scriptsize mutex +\\[-2pt]
     \scriptsize \code{SPI\_Read} 12\,B};
  \node[shade, below=2.9cm of ai] (af) {FT2232H};
  \draw[sig] (ai) -- node[font=\scriptsize, left, pos=0.4] {llave} (aw);
  \draw[sig] (ai) -- node[font=\scriptsize, right, pos=0.4] {llave} (ar);
  \draw[sig] (aw) -- node[font=\scriptsize, left] {1} (af);
  \draw[sig] (ar) -- node[font=\scriptsize, right] {2} (af);
  \node[note, below=0.25cm of af] (an) {dos transacciones por paso};
  \node[panel, fit=(ai)(aw)(ar)(af)(an), label={[font=\footnotesize\bfseries]
    above:Antes: bloques anteriores}] (pa) {};

  % ---------------- Ahora ----------------
  \begin{scope}[xshift=8.2cm]
  \node[box] (nw) {\code{escribirPic}\\[-1pt]\scriptsize deja el voltaje\\[-2pt]
    \scriptsize en memoria};
  \node[box, fill=black!12, below=1.0cm of nw, minimum width=3.4cm] (ni)
    {\code{initializePic}\\[-1pt]\scriptsize una \code{SPI\_ReadWrite} de 8\,B\\[-2pt]
     \scriptsize en \code{mdlUpdate}, sin mutex};
  \node[box, right=1.3cm of ni] (nr) {\code{leerPic}\\[-1pt]\scriptsize lee la\\[-2pt]
    \scriptsize posición};
  \node[shade, below=1.0cm of ni] (nf) {FT2232H};
  \draw[sig] (nw) -- node[font=\scriptsize, left] {voltaje} (ni);
  \draw[sig] (ni) -- node[font=\scriptsize, above] {posición} (nr);
  \draw[sig, <->] (ni) -- node[font=\scriptsize, left] {1} (nf);
  \node[note, below=0.25cm of nf] (nn) {una transferencia por paso};
  \node[panel, fit=(nw)(ni)(nr)(nf)(nn), label={[font=\footnotesize\bfseries]
    above:Ahora: bloques nuevos}] (pn) {};
  \end{scope}
\end{tikzpicture}
\caption{Operaciones sobre el FT2232H en cada paso, con los bloques anteriores
  y con los nuevos de los mismos nombres. En los nuevos, \code{escribirPic} y
  \code{leerPic} sólo leen y escriben memoria; \code{initializePic} es el único
  que usa el bus, así que no queda nada que proteger con un mutex.}
\label{fig:antes-ahora}
\end{figure}

Cómo trabajan los de tres bloques:

\begin{itemize}
\item \code{escribirPic} y \code{leerPic} ya no usan el bus. \code{escribirPic}
  deja el voltaje en la memoria compartida y \code{leerPic} lee la última
  posición recibida.
\item \code{initializePic} hace la única transferencia del paso, en su
  \code{mdlUpdate}. Simulink calcula las salidas de todos los bloques antes de
  cualquier actualización, así que el voltaje del paso ya está escrito y el
  orden entre bloques no importa.
\item La llave sigue siendo la dirección del estado compartido, ahora como
  \code{double}. Antes de usarla, cada bloque comprueba con
  \code{VirtualQuery} que la dirección sea válida y que el estado tenga una
  marca conocida. Una señal equivocada produce un mensaje de error en lugar de
  cerrar MATLAB.
\end{itemize}

Otros cambios respecto a los bloques anteriores:

\begin{itemize}
\item \textbf{Búsqueda de la interfaz.} Se prueba primero el canal cuyo
  número de serie termina en \code{A} y después los demás, con tramas que
  dejan la salida en 0\,V y no tocan el encoder; se usa el primero por el que
  el dsPIC responde con una trama válida (\code{protocolo/daq\_canal.h}).
  Funciona aunque libMPSSE numere las interfaces en otro orden y aunque el
  dsPIC se cablee a la interfaz~B. Si no responde por ninguna, el mensaje de
  error indica qué se leyó en cada una, por ejemplo ``interfaz A: solo 0xFF;
  interfaz B: solo 0xFF''.
\item \textbf{Verificación al arrancar.} El bloque pone en cero el encoder y
  comprueba que el dsPIC responda con el protocolo nuevo. Si no, detiene la
  simulación con un mensaje claro.
\item \textbf{Tiempo real.} Cada paso espera al reloj de la PC
  (\code{QueryPerformanceCounter}) y cuenta como atraso el que llega más de un
  periodo tarde.
\item \textbf{Errores con \code{ssSetErrorStatus}} en lugar de
  \code{mexErrMsgTxt}, y \code{mdlTerminate} tolera que Simulink lo llame sin
  haber llamado antes a \code{mdlStart}.
\item \textbf{Al terminar}, una trama con la salida deshabilitada (0\,V) y, en
  \code{initializePic}, un resumen en la ventana de comandos con pasos,
  errores y atrasos.
\end{itemize}

Para que MATLAB use los bloques nuevos y no los anteriores, la simulación se
ejecuta con \code{SIMULINK/DaqPic} como carpeta actual. El paquete para la
carpeta del asesor lo arma \code{herramientas/\allowbreak asesor/\allowbreak empaquetar.py}.

% ===========================================================================
\section{Herramientas de prueba sin Simulink}

Para separar las fallas del enlace de las de Simulink, el enlace se probó
primero sin MATLAB. Las herramientas están en \code{herramientas/} y comparten
el núcleo de la prueba (\code{herramientas/\allowbreak comun/\allowbreak enlace\_prueba.h}).

\begin{table}[H]
\centering
\small
\begin{tabularx}{\linewidth}{@{}lY@{}}
\toprule
Herramienta & Qué hace \\
\midrule
\code{protocolo/prueba\_protocolo.c} & Pruebas del formato de trama y del
  CRC, sin hardware, con un modelo del esclavo que reproduce el retraso de una
  trama. \\
\addlinespace
\code{prueba\_enlace} & Envía tramas al dsPIC y cuenta errores de encabezado,
  CRC y secuencia, los que reporta el dsPIC y la duración de cada
  transferencia. Compila en Windows con libMPSSE (\code{compilar.bat}) y en
  Linux con libftdi1. Con \code{-lazo} une MOSI y MISO dentro del FT2232H para
  probar la PC sin el dsPIC. \\
\addlinespace
\code{reporte/reporte.py} & Genera un PDF a partir del registro CSV de una
  prueba: errores, histograma de duraciones y mapa de bits erróneos. \\
\addlinespace
\code{visor\_enlace} & Interfaz gráfica en tiempo real que ejecuta el mismo
  lazo que un modelo de Simulink con \code{daqPic} (lazo abierto o PID), con
  osciloscopio, contadores y la tabla de tramas con error. Tiene un modo
  simulado sin hardware. \\
\addlinespace
\code{arnes\_simulink} & Ejecuta las S-Functions reales en Linux, sin MATLAB,
  llamándolas en el orden de Simulink contra la tarjeta. Revisa puertos,
  periodos, llaves inválidas, errores y la liberación del estado. \\
\addlinespace
\code{bateria\_pruebas.sh} & Ejecuta en secuencia el lazo interno y el dsPIC,
  lo más rápido posible y a 1\,kHz, y genera los reportes. \\
\bottomrule
\end{tabularx}
\caption{Herramientas de prueba.}
\label{tab:herramientas}
\end{table}

En Linux, la EEPROM del FT2232H está reprogramada con el identificador
\code{2099:0001}, y \code{herramientas/prueba\_enlace/70-daq-spi.rules} da
acceso al dispositivo sin \code{sudo}. El procedimiento completo está en
\code{PRUEBAS.md}.

% ===========================================================================
\section{Resultados}
\label{sec:resultados}

Con el firmware nuevo, el PLL y el cable de GND entre las tarjetas, el enlace
transmitió 100\,000 tramas sin errores, tanto a 1\,kHz como lo más rápido
posible. En Windows, \code{MergedPic.slx} se ejecutó sin modificaciones y sin
errores. Cada prueba está en \code{resultados/<fecha>}, con su registro CSV y
su reporte en PDF (tabla~\ref{tab:resultados}).

\begin{table}[H]
\centering
\small
\begin{tabularx}{\linewidth}{@{}lYY@{}}
\toprule
Fecha & Prueba & Resultado \\
\midrule
2026-10-07 & Windows 10, MATLAB R2025b: \code{MergedPic.slx} sin
  modificaciones con los bloques nuevos, 10\,001 pasos a 1\,kHz & 0 errores,
  1 atraso \\
\addlinespace
2026-10-07 & Windows: \code{DaqPicPrueba} y \code{DaqPicBloquesPrueba} con el
  modelo abierto en el editor & 4 a 10 errores por vigilancia: Simulink se
  detuvo cada 0.25\,s y algunas pausas pasaron de 50\,ms \\
\addlinespace
2026-10-06 & Rampa de $\pm1$\,V, 10\,000 tramas a 1\,kHz & 0 errores en el
  enlace; el DAC no respondió porque DIN y SYNC estaban en terminales
  equivocadas \\
\addlinespace
2026-10-06 & Con el cable de GND: 100\,000 tramas a 1\,kHz y 100\,000 a
  7408 tramas/s & 0 errores \\
\addlinespace
2026-10-06 & Sin cable de GND, ambas tarjetas en el mismo concentrador USB &
  0.025\,\% de las transferencias con error \\
\addlinespace
2026-10-05 & Con PLL, sin cable de GND, cada tarjeta en un puerto USB
  distinto & 0.7\,\% con error, por pulsos falsos en \code{CS} \\
\addlinespace
2026-10-05 & Sin PLL ($F_{cy} = 4$\,MHz), 100\,000 tramas a 1\,kHz & 83\,670
  errores de CRC: el dsPIC no alcanzaba a presentar el primer bit de cada
  byte \\
\bottomrule
\end{tabularx}
\caption{Resultados de las pruebas, de la más reciente a la más antigua.}
\label{tab:resultados}
\end{table}

El resultado de \code{MergedPic.slx} es el que responde a la pregunta en la
práctica: el mismo modelo, con los mismos tres bloques y en el mismo Windows,
dejó de fallar al cambiar el firmware y el uso del bus, no al quitar el mutex.

Dos hallazgos de estas pruebas no tienen que ver con el software:

\begin{itemize}
\item \textbf{El cable de GND entre la Curiosity Nano y la tarjeta del FT2232H
  es necesario.} Sin él aparecen pulsos falsos en \code{CS}, y el dsPIC rearma
  la respuesta a mitad de la trama.
\item \textbf{El cableado del DAC8554 estaba mal.} DIN estaba en RB7 en lugar
  de RB2 (SDO2) y SYNC en una terminal vecina de RB4. Por eso, en las pruebas
  con la salida habilitada, el motor no se movió y la posición se quedó en 0.
\end{itemize}

% ===========================================================================
\section{Pendientes}

El enlace y los bloques ya funcionan en Windows; falta la salida analógica y
el motor.

\begin{enumerate}
\item Corregir el cableado del DAC8554 (DIN a RB2, SYNC a RB4) y medir
  \code{VOUTA} con \code{prueba\_enlace -voltaje}.
\item Probar con el motor conectado: primero en lazo abierto y después con el
  PID de \code{crear\_modelo\_pid.m}.
\item Comprobar en la interfaz gráfica de Simulink, no sólo con
  \code{matlab -batch}, si hay pausas de más de 50\,ms con el modelo abierto.
\item Ejecutar \code{MergedPic.slx} en la PC del laboratorio con el paquete de
  \code{herramientas/asesor}.
\end{enumerate}

\end{document}
